\section{Introduction}

Hygiene-augmented vulnerability prioritization is a deployment, not a
demonstration. The HygienePrio scorer~\cite{paper4hygieneprio} adds
approximately 31~pp of Precision@50 over EPSS-only ranking at a single
window, and Paper~5~\cite{paper5temporal} showed the per-pair advantage
persists across six rolling maintenance windows. But Paper~5 also
identified that the \textit{absolute} P@50 produced by the
fixed-Paper-4-weights scorer is regime-dependent, and an
offline-peek calibration ablation could add up to $+21.3$~pp at the
W2 window of the canonical $(K=50, \lambda=3)$ cell. The peek
procedure is not deployable. The question this paper answers is:
can any deployable recalibration procedure recover the peek
ceiling, and if so, in which capacity regime?

We attack the question in pre-registered increments. Six experiments,
each addressing the failure mode the previous one exposes, are
locked-and-evaluated in sequence:

\begin{itemize}
\item \textbf{E1.} Lag-1 rolling-history grid search --- the simplest
deployable substitute for the offline peek.
\item \textbf{E2.} Multi-window smoothing (EWMA-3, trailing-mean-3)
to address E1's failure mode.
\item \textbf{E3.} Magnitude-based change-point detection with a
single threshold $\tau$ to switch between lag-1 and fixed.
\item \textbf{E4.} Threshold-sensitivity sweep over $\tau$ to map
the H1/H2 Pareto frontier of E3.
\item \textbf{E5.} Capacity-aware $\tau$ vector $\tau_K$, locked from
inspection of E4's calibration-seed log.
\item \textbf{E6.} One-sided CUSUM with pre-registered $(k, h)$ to
address E5's structural-collapse property.
\end{itemize}

Each experiment is governed by its own pre-registration document
(archived in \texttt{paper7/experiments/0X\_*/PROTOCOL.md}); 27
hypotheses are locked across the six experiments and reported
honestly here. The full hypothesis-outcome ledger is in
Table~\ref{tab:hypothesis_ledger}.

\subsection{Contributions}

\begin{itemize}
\item \textbf{C1 --- A six-experiment pre-registered map of the
deployable-calibration design space.} 27 locked hypotheses, 6
sequential frozen evaluations on the EEHDA fleet, 19{,}500
result rows total. Each experiment's pre-registration was locked
before its evaluation-seed P@50 was inspected; each rejection is
reported per the protocol's stop rules.

\item \textbf{C2 --- Lag-1 works at moderate capacity but harms at
high capacity (E1).} Cell-mean recovery ratio of the offline ceiling
is $\rho \in \{1.04, 0.99\}$ at $K \in \{50, 100\}$ and
$\rho = -0.66$ at $K = 200$. The capacity-conditional hazard is the
substantive finding that drives the rest of the program.

\item \textbf{C3 --- Multi-window smoothing falsifies the
variance-reduction prior (E2).} EWMA-3 and trailing-mean-3
\textit{worsen} the K=200 hazard ($\rho = -0.94$ for EWMA-3 vs
$\rho = -0.66$ for lag-1) because older calibration windows are
mis-aligned with the current window in a fast-shift regime.
``More history reduces variance and helps'' is wrong in this regime.

\item \textbf{C4 --- Single-$\tau$ magnitude detector cannot reach
joint feasibility (E3, E4).} A single threshold $\tau$ that avoids
the K=200 hazard also fires at K=50 where lag-1 helps; the swept
Pareto curve sits entirely outside the joint feasibility region
defined by Papers~5/10's two tolerances.

\item \textbf{C5 --- Capacity-aware $\tau$ reaches feasibility but
collapses to the static gate (E5).} A per-K $\tau_K$ vector with
$\tau_{200} = 0.02$ satisfies all three pre-registered tolerances,
but its K=200 cell mean equals the static gate's cell mean exactly
($0.2664$ to four decimal places) because the small $\tau_{200}$
fires at every post-W1 K=200 window. The per-window detector
contributes nothing beyond the static gate at the high-capacity
extreme.

\item \textbf{C6 --- CUSUM is the first detector to beat the static
gate (E6).} A one-sided CUSUM with $k = 0.04$, $h = 0.10$ filters
single-window noise excursions (notably the $K=200, w=4$ excursion
where $|\Delta| = 0.020 < k$) and uses the lag-1 fit at those
windows. K=200 cell-mean P@50 = $0.2754$, beating the static gate's
$0.2664$ by $+0.9$~pp. This is \textit{the only operationally
meaningful gain over the static gate in the six-experiment sequence};
the H2.E6 binary test rejects by a fine 0.001~pp margin, which we
report honestly. CUSUM with a single $(k, h)$ pair over-fires at
K=50 ($-3.9$~pp cell-mean cost), the natural target of a
capacity-aware CUSUM successor.

\item \textbf{C7 --- Capacity-conditional deployable recipe.}
At $K \leq 100$ deploy lag-1 (E1, $\rho \approx 1.0$). At $K = 200$
deploy CUSUM (E6, $+0.9$~pp over the static gate). For variable
capacity across deployment cycles, a per-K $(k_K, h_K)$ CUSUM
vector is the natural successor and is left as future work under
pre-registration discipline.
\end{itemize}

\subsection{Relationship to Prior Papers}

This paper sits between Paper~5~\cite{paper5temporal} (temporal
stability at fixed weights), Paper~6~\cite{paper6capacity} (capacity
sweep that identified the K=200 collapse), and the methodological
self-trajectory paper~\cite{paper9selftraj} (which addresses the
selection-policy coupling threat that all six experiments here
inherit). The present paper takes the K=200 absolute-floor problem
that Paper~6 identified and asks whether deployable online
recalibration can recover the offline ceiling. The honest answer
unifying all six experiments is: yes at moderate capacity; only
CUSUM-class detectors at high capacity, and only marginally.
